% Paul E. West and Sean T. Hayes

\documentclass[14pt,aspectratio=1610]{beamer}
% \documentclass[14pt,aspectratio=1610, handout]{beamer}
\usepackage{csu-theme}

\title{\Huge{}{\ttfamily\bfseries C++} Pointers and Testing}
\subtitle{CSCI 315: Data Structure Analysis}
\institute{Charleston Southern University}

%% Comment out date to use default (today)
% \date{January 23, 2025}
\subject{Software Programming}
%\keywords{}

\begin{document}

\begin{frame}
  \titlepage{}
\end{frame}

\section{Overview}


\begin{frame}{Overview}
\begin{itemize}[<+->]
\item Objects in C++ share many similarities with Java, \\
  but are considered more powerful:
\begin{itemize}
\item They can be created on the stack.
\item C++ allows multiple inheritance.
\begin{itemize}
\item no need for the interface keyword.
\end{itemize}
\item Templates act like generics, except they are a compile-time construct, not run time.
\end{itemize}
\item Remember, C++ is compiled to machine code; Java compiles to byte code!
\end{itemize}
\end{frame}


\section{Review}

\begin{frame}[fragile]{Pointers to Stack Memory}
\renewcommand{\arraystretch}{1.5}
\begin{tabularx}{1\textwidth}{>{\columncolor{codeBackground}}l >{\small}X}
  \lstinline{int num = 32;} & Declare an \lstinline{int} variable on the stack and set to \lstinline{32}.
    \onslide<2->\\
  \lstinline{int *ptr;}   & Declare a variable \lstinline{ptr} to store the memory address of where an \lstinline{int} resides.
    \onslide<3->\\
  \lstinline{ptr = \&num;}   & \lstinline{\&num} gets the memory address of \lstinline{num}.\newline
  Stores the memory address of \lstinline{num} in \lstinline{ptr}.
    \onslide<4->\\
  \lstinline{*ptr = 64;}  & Go to where \lstinline{ptr} is pointing (\lstinline{*}) and store \lstinline{64} there.\newline
  Since \lstinline{ptr} is pointing to \lstinline{num}, \lstinline{num} now has \lstinline{64}.\onslide\\
\end{tabularx}

% \begin{itemize}
% \item \lstinline{int *a;}
% \begin{itemize}
% \item
% \end{itemize}
% \item \lstinline{int b = 32;} \\
% \lstinline{a = \&b;}
% \item
% \item
% \item \lstinline{*a = 64};
% \item
% \item
% \end{itemize}
\end{frame}

%\setbeamercovered{invisible}
\begin{frame}{Pointer to Heap Memory}
\renewcommand{\arraystretch}{1.5}
\begin{tabularx}{1\textwidth}{>{\columncolor{codeBackground}}l >{\small}X}
  \lstinline{int *a;}
    & Create a variable \lstinline{a} on the stack that will store the memory address of an integer.\onslide<2->\\
  \lstinline{a = new int;}
    & Allocate an \lstinline{int} in memory and store the memory address of that \lstinline{int} in \lstinline{a}.\onslide<3->\\
  \lstinline{*a = 5;}
    & Manipulate the new \lstinline{int} just like before.\onslide<4->\\
  \lstinline{delete a;}
    & Because the integer is on the heap, you must de-allocate!\onslide<5->\\
  \lstinline{a = nullptr;}
    & After \lstinline{delete}, reassign the pointer to \lstinline{nullptr} (best practice).\onslide\\
\end{tabularx}
\end{frame}

\begin{frame}[fragile]{Arrays with Pointer Arithmetic}
\renewcommand{\arraystretch}{1.5}
\begin{tabularx}{1\textwidth}{>{\columncolor{codeBackground}}l X}
  \lstinline{int ary[100];}
    & Create an array of 100 integers\newline on the stack.\\
  \lstinline{int *pAry = ary;}
    & Create a pointer to the array.\newline Note: \lstinline{\&} was not used.\\
\end{tabularx}\\[1em]
\pause
The following four statements are equivalent.\\
\renewcommand{\arraystretch}{0.5}
\begin{tabularx}{1\textwidth}{r l r}
  \cellcolor{codeBackground}\lstinline{ary[10] = 50;}
    && \cellcolor{codeBackground}\lstinline{*(ary + 10) = 50;}\\
    &&\\
     \cellcolor{codeBackground}\lstinline{pAry[10] = 50;}
    && \cellcolor{codeBackground}\lstinline{*(pAry + 10) = 50;}\\
\end{tabularx}\\[1em]

\end{frame}

\begin{frame}[fragile]{Arrays}
\begin{itemize}
\item In Java, arrays are actually objects that store length and other information.
\item In C++, an array is a sequence list of data. (Nothing else is stored.)

\item These are valid declarations in C++, but not in Java. Why?
\begin{lstlisting}
int x[20];    // Create 20 integers
Button b[20]; // Create 20 Buttons
\end{lstlisting}
\end{itemize}
\end{frame}

\begin{frame}{Pointers and Arrays}
\begin{itemize}
\item In C++, pointers and arrays have a close relationship.
\item Instead of \lstinline{int x[20];}\\
we can can write \\
\lstinline{int *x = new int[20];} \\
to allow for dynamic allocation.
\begin{itemize}
\item Usage of the array (e.g., \lstinline{x[3] = 5;}) identical in both cases
\item To deallocate, use \lstinline{delete[] x;}
\end{itemize}
\end{itemize}
\end{frame}

\section*{Constants}
\begin{frame}{Pointers and Const-Correctness}
  A good rule of thumb is to use \lstinline{const} wherever possible.
  Pointers have two types of const-ness.
  \begin{itemize}[<+->]
    \item A pointer to constant data\\
      \lstinline{const int * ptr;}\\
      \lstinline{int const * ptr;}\\
      You cannot change the value at the address this pointer is pointing to.
    \item A pointer with a const memory address\\
      \lstinline{int * const ptr;}\\
      You cannot change what this pointer points to.
  \end{itemize}
\end{frame}

\begin{frame}{Pointers and Const-Correctness}
  A good rule of thumb is to use \lstinline{const} wherever possible.
  Pointers have two types of const-ness.
  \begin{itemize}[<+->]
    \item A const pointer to const data\\
      \lstinline{const int * const ptr;}\\
      \lstinline{int const * const ptr;}
  \end{itemize}
\end{frame}

\section{Object Pointers}

\begin{frame}{In C++, a pointer can point to anything.}
  \begin{itemize}
  \item Primitives
  \item Arrays
  \item Other pointers
  \item Objects
  \item Functions
  \item Something nonspecific
  \end{itemize}
\end{frame}

\begin{frame}[fragile]{Casting}

C++ is incredibly powerful, because of casting.\\
\begin{lstlisting}
int num = 300;
int* const pInt = &num;
float* const pFloat = reinterpret_cast<float*>(pInt);
*pFloat = 3.14159; // undefined behavior
std::cout << "*pFloat = " << *pFloat;
std::cout << "\n    num = " << num << '\n';
\end{lstlisting}

\begin{itemize}
\item
  Yes, that is ``legal.''
\item 
  The value of \lstinline{num} is \underline{\alt<2>{undefined}{\phantom{undefined}}}!
\end{itemize}
\end{frame}

\begin{frame}[fragile]{Pointers to Objects}
\begin{columns}
\column{0.55\textwidth}
Pointers to instances of \lstinline{struct}s and \lstinline{class}es
work the same as with primitive types (e.g., \lstinline{int}).
\column{0.35\textwidth}
\begin{lstlisting}[basicstyle=\footnotesize\ttfamily, xleftmargin=0pt]
struct Dog {
  std::string name;
  std::string breed;
  int yearsOld;
};
\end{lstlisting}
\end{columns}
\begin{minipage}{\dimexpr\textwidth-2\fboxsep-2\fboxrule-60pt} 
\begin{lstlisting}[basicstyle=\footnotesize\ttfamily]
int main(int argc, char *argv[]) {
  Dog myPet {"Bella", "German Shepherd", 7};
  Dog *pFriend = new Dog {"Rocky", "Poodle", 2};

  std::cout << myPet.name << std::endl;
  std::cout << (*pFriend).name << std::endl;

  // using member-access arrow
  std::cout << pFriend->name << std::endl;

  return 0;
}
\end{lstlisting}
\end{minipage} 

\end{frame}

\begin{frame}{Pointers to Objects}
\begin{itemize}
\item In Java, a pointer is called a reference to an object:\\
\lstinline[language=Java]{String str = new String("Hello!");}
\item In C++, a pointer stores the memory location of something else. For objects:\\
\lstinline{std::string *str = new std::string("Hello!");}

\end{itemize}
\end{frame}

\begin{frame}{Pointers to Objects: A Simple Explanation}
\begin{itemize}\large
\item In Java, pointers only exist for objects.
\item In C++, a pointer may point to anything.
\end{itemize}
\end{frame}

\section{Function Pointers}
\begin{frame}{Function Pointers}
\begin{itemize}[<+->]
  \item
    A \emph{pointer} is a variable that stores the address of something.
  \item
    A \emph{function pointer} stores the address of a function.
  \item
    The function can be called by dereferencing the function pointer.
  \item
    Function pointers can be passed as parameters to other functions.
  \item
    Functions can return function pointers.
  \item
    Commonly used in \href{https://en.wikipedia.org/wiki/Callback_(computer_programming)}{callbacks}.
\end{itemize}
\end{frame}

\begin{frame}[fragile]{Function Pointers}
\begin{lstlisting}[basicstyle=\scriptsize\ttfamily]
#include <iostream>

int foo() { return 4; }

int main(int argc, char *argv[])
{
  // Display functions memory address.
  // Cast to void pointer to ensure the address is properly displayed.
  std::cout << reinterpret_cast<void*>(main) << std::endl;
  std::cout << reinterpret_cast<void*>(foo) << std::endl;

  // Create a pointer that points to foo()
  int (*fcnPtr)() = foo;

  // Call foo() using pointer
  std::cout << "The return value of foo() is "
    << (*fcnPtr)()
    << std::endl;

  return 0;
}
\end{lstlisting}
\end{frame}

\begin{frame}{There is more}
Pointers can point to many things.
\begin{itemize}
\item Primitives
\item Arrays
\item Other pointers
\item Objects
\item Functions
\item Something nonspecific (\lstinline{void *}).\\
  Void pointers are more useful in C but should be avoided in C++.
\end{itemize}
\end{frame}


\section*{Lab 4}
\begin{frame}{Lab 4: Pointers}
\centering\Large
Let's talk about Lab 4, which is more practice with pointers.
\end{frame}


\section*{Ethics Essay}

\begin{frame}{Ethics Essay: THERAC-25}
\centering\Large
\href{https://github.com/csu-cs/CSCI-315-2025-Fall/blob/master/ethics/README.md}{Let's go over your ethics essay today!}
\end{frame}

\section{Test Your Code}

\begin{frame}{Testing}
\textbf{\large\textit{Testing will make you more productive.}}\\[1em]
Meaning, you will spend less time working on problems because you will\textellipsis
\begin{itemize}
\item
  increase your chances of getting it right the first time AND
\item
  rarely make the same mistake twice.
\end{itemize}
\end{frame}

\begin{frame}{Warning}
\begin{itemize}[<+->]
\item A computer program \textit{cannot} sufficiently generate test cases for your code.
\item \href{https://en.wikipedia.org/wiki/Halting_problem}{The Halting Problem}: Alan Turing proved that it is impossible for a computer to determine if any program halts (finishes).\\
(\href{https://youtu.be/92WHN-pAFCs?t=14}{see Explanation on YouTube.})
\item If we don't know if it halts, how can we test?
\item This is your job, and an advantage you have over computation.
\end{itemize}
\end{frame}

\begin{frame}{Unit Test}
\begin{itemize}
\item
  For this course, unit testing is sufficient.
\item
  There are many other forms of testing, but this course focuses on (relatively) small units of code.
\item
  There are numerous unit-testing tools, but no \textit{de facto} standard for C++.
  \begin{itemize}
  \item Unlike Java's JUnit
  \end{itemize}
\item
  Popular testing frameworks for C++ include \href{https://google.github.io/googletest/}{GoogleTest}, \href{https://www.boost.org/doc/libs/1_81_0/libs/test/doc/html/index.html}{Boost.Test}, \href{CppUnit}{CppUnit}, \href{https://cxxtest.com/}{CxxTest}, and \href{https://github.com/catchorg/Catch2}{Catch2}.
\end{itemize}
\end{frame}

\begin{frame}{Doctest -- Unit Testing for C++}
\begin{itemize}
\item
  We will use \href{https://github.com/doctest/doctest\#readme}{Doctest} for this course.
  \begin{itemize}
  \item
    Doctest is simple, yet has a rich feature set.
  \end{itemize}
\item
  You only need to place the one file, \href{https://github.com/doctest/doctest/blob/master/doctest/doctest.h?raw=true}{doctest.h}, in your test repository.
\item
  I have provided a full template in class-code.\\Let's take a look\textellipsis
\end{itemize}
\end{frame}

\section*{Lab 5}

\begin{frame}{Auto-Grader}
  \begin{itemize}
  \item An auto-grader will provide a grade for most of the labs (starting in Lab 5).
  \begin{itemize}
    \item It will not say what test cases you failed.
    \item It may not mention if there is a timeout.
    \item It may not say if there are compilation errors.
    \item It \textit{may} tell you which part of the lab you got right or wrong.
  \end{itemize}
  \item Test cases from the auto-grader will \textbf{NOT} be disclose.  You \textbf{must} learn how to write your own test cases!
  \item The grade from the auto-grader should be pushed to your repository shortly after you push your solution.
  \end{itemize}
\end{frame}

\begin{frame}{Lab 5: A Review of C++ Arrays}
\centering\Large
Let's talk about Lab 5.
\end{frame}

\end{document}
